\documentclass[11pt]{article}
\usepackage[a4paper,margin=25mm]{geometry}
\usepackage{amsmath,amssymb,amsthm,hyperref}
\newtheorem{theorem}{Theorem}\newtheorem{lemma}{Lemma}
\newtheorem{corollary}{Corollary}
\title{An elementary quadratic bound for the last atom}
\author{Complete elementary proof; three independent mathematical reviews}
\date{30 September 2026}
\begin{document}\maketitle
Let $w_q>0$, $1\le q\le K$, sum to one, and let
$p_i(q)\in[0,1]$, $1\le i\le m$, be nondecreasing in $q$. Suppose
\begin{equation}\label{mom}
 \sum_qw_q\prod_{i\in S}p_i(q)=\frac1{|S|+1}
 \qquad(S\subseteq\{1,\ldots,m\}).
\end{equation}
The empty-set equation is the normalization. No equal weights or
rationality are assumed. Write $a_i=p_i(K)$.

\begin{theorem}\label{main}
For every $m\ge1$,
\[
 a_i\ge1-\frac{\pi^2}{(m+1)^2}\quad(1\le i\le m),
 \qquad w_K\le\frac{200}{m^2}.
\]
The second bound also holds for the total weight of all rows identical
to the last vector. Consequently an equal-weight representation has
$K\ge m^2/200$.
\end{theorem}

The proof uses only squarefree moment exactness, monotonicity, and
ordinary Gaussian and Radau quadrature. In particular it gives a short
independent proof of the exact quadratic per-die lower bound, without
the endpoint localization estimates developed earlier in this project.

\begin{lemma}[Paired thresholds]\label{pair}
Complete the coordinatewise ordered row vectors to a continuous
nondecreasing path by adjoining the vectors $0$ and $1$ and linearly
interpolating successive vectors. For distinct indices $j,k$ and
$t\in(0,1)$, stop at the first time either coordinate reaches $t$.
The resulting threshold values $\alpha,\beta$ satisfy
\[
 \{\alpha,\beta\}=\{t,b\},\quad 0\le b\le t,
 \qquad (p_j(q)-\alpha)(p_k(q)-\beta)\ge0
 \quad\hbox{for every original row }q.
\]
Products of these factors on disjoint pairs remain nonnegative.
\end{lemma}
\begin{proof}
Both coordinates start at zero and end at one. At their first crossing
of $t$, one equals $t$ and the other is at most $t$. Every original row
before the stopping point is coordinatewise below the threshold vector,
and every row after it is coordinatewise above it. Thus the two
differences have the same weak sign. The added path points have no
probability weight and change none of the moment equations.
\end{proof}

For every multi-affine polynomial $F$, equation~\eqref{mom} implies
the exact diagonal identity
\begin{equation}\label{diagonal}
 \mathbb E_w F(p_1,\ldots,p_m)=\int_0^1F(t,\ldots,t)\,dt.
\end{equation}
This is simply linearity applied to its squarefree monomials.

\begin{lemma}[Quadrature facts]\label{quad}
The $D$-node Gaussian rule for Lebesgue measure on $[0,1]$ has positive
weights, is exact through degree $2D-1$, and its largest node $\xi_D$
satisfies
\begin{equation}\label{gaussnode}
 1-\xi_D\le\frac{\pi^2}{4(D+1)^2}.
\end{equation}
The $(e+1)$-node right Radau rule has nodes
$t_1<\cdots<t_e<1$ and $1$, positive weights, and exactness through
degree $2e$. Its endpoint weight and interior-node sum are
\begin{equation}\label{radau}
 \omega_1=\frac1{(e+1)^2},\qquad
 \sum_{\ell=1}^e\frac1{1-t_\ell}=\frac{e(e+2)}2.
\end{equation}
\end{lemma}
\begin{proof}
Let $P_r$ be the usual Legendre polynomial with $P_r(1)=1$.
Orthogonality yields the Gaussian rule at the shifted roots of $P_D$:
division by the node polynomial gives exactness, and applying the rule
to squared cardinal polynomials gives positive weights.
The three-term Legendre recurrence shows that its unshifted roots are
the eigenvalues of the symmetric $D\times D$ Jacobi matrix with zero
diagonal and off-diagonal entries
$k/\sqrt{(2k-1)(2k+1)}\ge1/2$. The positive sine eigenvector of the
matrix with all off-diagonals $1/2$ gives, by the Rayleigh principle,
largest eigenvalue at least $\cos(\pi/(D+1))$. Shifting to $[0,1]$
and using $1-\cos x\le x^2/2$ proves~\eqref{gaussnode}.

For Radau put
\[
 J_e(t)=\sum_{r=0}^e(2r+1)P_r(2t-1).
\]
This is the reproducing kernel at $1$ for degree-$e$ polynomials:
$\int_0^1J_e(t)f(t)dt=f(1)$. Hence $J_e$ is orthogonal to degree
at most $e-1$ for the positive weight $1-t$, and has $e$ distinct
interior roots. Interpolation at these roots and $1$ integrates every
polynomial of degree at most $2e$: the division remainder has degree
at most $e$, and the quotient term is $(t-1)J_e(t)q(t)$ with
$\deg q\le e-1$. Interior weights are positive by evaluating the
rule on
$(1-t)(J_e(t)/(t-t_\ell))^2$. At the endpoint, exactness on $J_e^2$
and the reproducing identity give
\[
 \omega_1 J_e(1)^2=\int_0^1J_e(t)^2dt=J_e(1)=(e+1)^2.
\]
Finally the elementary endpoint derivative formula for shifted Legendre
polynomials is $(P_r(2t-1))'|_{t=1}=r(r+1)$. Therefore
\[
 \frac{J_e'(1)}{J_e(1)}
 =\frac{\sum_{r=0}^e(2r+1)r(r+1)}{(e+1)^2}
 =\frac{e(e+2)}2.
\]
The left side is the logarithmic derivative over the roots, proving
the second identity in~\eqref{radau}.
\end{proof}

\begin{proof}[Proof of Theorem~\ref{main}]
Fix $i$ and put $d=\lfloor(m-1)/2\rfloor$, $D=d+1$.
Choose $d$ disjoint pairs among the other $m-1$ columns. For pair
$\ell$, apply Lemma~\ref{pair} at the $\ell$th node of the
$D$-node Gaussian rule, omitting its largest node $\xi_D$.
Let $T(p)$ be the product of the $d$ paired factors and put
\[
 P(t)=T(t,\ldots,t)=\prod_{\ell=1}^d(t-\xi_\ell)(t-b_\ell),
 \qquad b_\ell\le\xi_\ell<\xi_D.
\]
Both $T(p(q))$ and $a_i-p_i(q)$ are nonnegative. They use disjoint
columns, so their product is multi-affine, of degree $2d+1$. By
\eqref{diagonal} and Gaussian exactness,
\[
 0\le\mathbb E_w[(a_i-p_i)T(p)]
 =\int_0^1(a_i-t)P(t)dt
 =\omega_D(a_i-\xi_D)P(\xi_D).
\]
Here $\omega_D>0$ and $P(\xi_D)>0$, including the empty-product
case $d=0$. Thus $a_i\ge\xi_D$ for every $i$. Since
$D+1\ge(m+1)/2$, equation~\eqref{gaussnode} gives
$1-a_i\le\delta:=\pi^2/(m+1)^2$.

For $m\ge10$, put $e=\lfloor m/10\rfloor$ and choose $e$ disjoint
pairs of columns. Apply the threshold lemma to these pairs at the
$e$ interior nodes of the right Radau rule. Again write $T$ for the
nonnegative multi-affine product and
\[
 P(t)=\prod_{\ell=1}^e(t-t_\ell)(t-b_\ell),\qquad b_\ell\le t_\ell.
\]
The degree is $2e\le m$. Radau exactness and~\eqref{diagonal} give
\begin{equation}\label{radauidentity}
 \mathbb E_w T(p)=\frac{P(1)}{(e+1)^2}.
\end{equation}
Set $S=e(e+2)/2$. As $e\le m/10$ and $e+2\le3m/10$,
\[
 \delta S\le\frac{3\pi^2}{200}<\frac14.
\]
In particular all threshold roots are strictly below every $a_i$.
At the last row each paired factor, divided by its value at the
all-one vector, is at least
$(1-\delta/(1-t_\ell))^2$. The elementary inequality
$\prod_j(1-z_j)\ge1-\sum_jz_j$ for $0\le z_j\le1$ therefore gives
\[
 \frac{T(a)}{P(1)}
 \ge\prod_{\ell=1}^e
       \left(1-\frac{\delta}{1-t_\ell}\right)^2
 \ge1-2\delta S>\frac12.
\]
Consequently~\eqref{radauidentity} implies
$w_K\le2/(e+1)^2\le200/m^2$. Identical terminal rows contribute
the same value $T(a)$, giving the aggregate version. For $m<10$ the
asserted bound follows from $w_K\le1$.
\end{proof}

\begin{corollary}[Fair dice]
Every die in a permutation-fair family of $n\ge2$ finite uniform dice
has at least $(n-1)^2/200$ faces.
\end{corollary}
\begin{proof}
Condition on the sorted faces of that die. The comparison probabilities
of the other $m=n-1$ dice are monotone and satisfy~\eqref{mom}.
Each conditioning row has weight the reciprocal of its face count.
Apply Theorem~\ref{main}.
\end{proof}
Only the mixed moment conditions~\eqref{mom} were used. Thus full
permutation fairness can be weakened to uniform winning probability
for this distinguished die on every subset containing it.
\end{document}
